\section{Explicit low graph residues: genuine Tor classes and the mixed obstruction}
\label{sec:136-low-residue-formulas}

\paragraph{Archive and scope.}
These supplementary calculations concern cote136, physical
pp.~117,119,121, footers116,118,120, scans117,119,121.
They explain two pieces of ARCH-136-Q3-LOW-FILTRATION.
The first gives actual common graph primitives for a genuine
integer Tor class. The second computes the entire finite mixed
piece and shows that its residue is determined by the old maps.
Neither calculation asserts an exhaustive Tor generation theorem,
preservation of all graph residues, or a complete common source.

Write \(S=\Gamma_{\mathbb Z}(T)\), \(J=\mathbb Q[T]/S\), and
\(G=\partial\Delta[2]\), with its usual boundary orientation.
All source maps below are actual strict maps into admissible
targets \(C_j\) having the literal common first skeleton \(F\).
They agree on that entire skeleton. The fixed comparison and row
bounds are chosen so that the target vertical groups vanish in
positive degrees through three and the graph comparison gives
\[
 H^0(F^{w,*}(G))=S^w,\qquad H^1(F^{w,*}(G))=S^w,\qquad
 H^q(F^{w,*}(G))=0\quad(q\geq2)
                                                        \tag{136.LGF1}
\]
in the degrees used. Let \(\operatorname{per}_G\) denote the
oriented identification of the degree-one group
with \(S^w\).

For a source three-cycle \(\xi\) with \(dB=L\xi\), target
primitive pairs satisfy
\[
 dP_j=f_j\xi,\qquad dQ_j=f_jB-LP_j.
                                                        \tag{136.LGF2}
\]
When the \(P_j\) have common graph restrictions, its graph
difference is
\[
 \left[\frac{\operatorname{per}_G((Q_j-Q_0)|_G)}{L}\right]
       \in J^w.                                        \tag{136.LGF3}
\]
If source second cohomology is integer torsion, this is the
universal linking class independent of representative and order
relation: changing the relation produces the ordinary graph-period
map from that second cohomology to the torsion-free \(S^w\),
which is zero. This is the exact mechanism already proved in
ARCH-136-TWO-TRUNCATED-LINKING-REPAIR and
ARCH-136-TORSION-GRAPH-RESIDUE. For the first calculation one may
instead simply retain the particular displayed Moore pair.

\subsection{Actual primitives for a genuine Tor class}

\paragraph{ARCH-136-TOR-GRAPH-PRIMITIVES.}
Let \(\kappa_i\) be source two-cycles and let \(b_i\) be
one-forms with
\[
 db_i=N_i\kappa_i\quad(i=1,2),\qquad N_i\geq1.
\]
Assume each individual weight is in the fixed common degree-two
primitive range. Put
\[
 \begin{gathered}
 g=\gcd(N_1,N_2),\quad \alpha=N_1/g,\quad \beta=N_2/g,\quad P=N_1N_2/g,\\
 \xi=\beta b_1\kappa_2-\alpha\kappa_1b_2.
\end{gathered}
                                                        \tag{136.LGF4}
\]
Then \(d\xi=0\), \(g\xi=d(-b_1b_2)\), and this actual
Moore pair has target primitives with common graph restrictions.
In particular its graph difference is zero.

\begin{proof}
Use the common degree-two primitive lemma to choose, in each target,
\[
 da_i=f\kappa_i,\qquad dr_i=N_i a_i-fb_i,
                 \qquad |a_i|=1,\quad |r_i|=0,          \tag{136.LGF5}
\]
with all graph restrictions common across targets. Define
\[
 A=-P a_1a_2-\beta r_1\,da_2+\alpha a_1\,dr_2,
 \qquad Q=r_1\,fb_2.                                   \tag{136.LGF6}
\]
The first assertion follows by differentiating (136.LGF4), using
\(N_1\beta=N_2\alpha=P\). For the primitives, the cochain
signs give
\[
\begin{aligned}
 dA
 &=P(a_1\,da_2-da_1\,a_2)
            -\beta\,dr_1\,da_2+\alpha\,da_1\,dr_2\\
 &=\beta(fb_1)\,da_2-\alpha\,da_1(fb_2)=f\xi.
\end{aligned}                                         \tag{136.LGF7}
\]
Expanding \(fb_i=N_i a_i-dr_i\), the same identities give
\[
 f(-b_1b_2)-gA
    =N_2\,dr_1\,a_2+N_2r_1\,da_2-dr_1\,dr_2
    =d(r_1\,fb_2).                                    \tag{136.LGF8}
\]
Thus (136.LGF6) is the required pair. Every expression is a
polynomial in elements whose graph values are common, so its graph
values are common too. No coefficient is divided in the source or
targets. The separate weight guards in (136.LGF5) are essential;
a bound only on the sum of the two weights would not justify them.
\end{proof}

This supplies an explicit version of the free-unit-disk argument in
ARCH-136-Q3-ISOLATED-TOR. It does not cancel a positive coefficient
from the coefficient-square construction of that result.

\subsection{The exact residue of a finite mixed class}

Let \(A\) be a two-truncated source with \(A_0=S\),
\(A_1=F_1\), \(H^0(A_2)=S\), \(H^1(A_2)=0\), and
rational coefficient cohomology. Suppose its old degree-three
linking differences vanish in the fixed output tail. Consider an
actual free full-boundary degree-three Tate stage with pointwise
presentation
\[
 \begin{gathered}
 E_2=A_2\langle t_i,c_i\rangle,\qquad |t_i|=1,\quad |c_i|=2,\\
 dc_i=z_i,\quad dt_i=N_i c_i-b_i,\quad db_i=N_i z_i,\quad N_i\geq1.
 \end{gathered}
                                                        \tag{136.LGF9}
\]
Its complete boundary values are in the old first skeleton.
\emph{Assume actual extensions \(f_j:E\to(C_j)_{\leq2}\)
are given and agree on the entire first skeleton of \(E\).}
Only finitely many indices occur in the following expressions.

Take old elements \(x_i,y_i,a\), of form degrees \(2,1,3\),
such that
\[
 dx_i=0,\qquad dy_i=N_i x_i,\qquad
 da=\sum_i(x_i b_i-y_i z_i).                            \tag{136.LGF10}
\]
All summands have the same total weight \(w\).
Each individual weight of \(x_i,y_i\) is required to be in the
common degree-two primitive range; the total weight \(w\) is in
the old graph-compatibility and target comparison ranges.
Choose a common multiple \(L\) of the \(N_i\), and write
\[
 \begin{gathered}
 \ell_i=L/N_i,\qquad X=a+\sum_i(x_i t_i-y_i c_i),\\
 h=L a-\sum_i\ell_i y_i b_i.
\end{gathered}
                                                        \tag{136.LGF11}
\]

\paragraph{ARCH-136-MIXED-H2-RESIDUE.}
The elements \(X,h\) are three-cycles and
\[
 LX=d\left(\sum_i\ell_i y_i t_i\right)+h.                \tag{136.LGF12}
\]
The universal graph difference of \(X\) is the explicit old-data
secondary period in (136.LGF19) below. In particular, changing the
new target images of the \(c_i,t_i\), while retaining the old
maps and the common first-skeleton values, cannot change this
mixed residue. The old graph-compatibility hypothesis alone does
not supply the vanishing of the period in that formula.

\begin{proof}
Differentiating (136.LGF10)--(136.LGF11) proves \(dX=dh=0\)
and (136.LGF12). Rational coefficient cohomology gives an
integer \(M>0\) and an old two-form \(H\) with
\(dH=Mh\). The free extension (136.LGF9) is rationally a
unit-disk extension, by \(c_i'=c_i-b_i/N_i\). Thus its
second cohomology is integer torsion, as required for the
universal linking interpretation (136.LGF3).

Choose common graph primitive pairs for the old two-cycles:
\[
 du_i=f_jx_i,\qquad dv_i=N_i u_i-f_jy_i,
                         \qquad |u_i|=1,\quad |v_i|=0. \tag{136.LGF13}
\]
Target-dependent subscripts are suppressed inside formulas.
The old-data three-form
\[
 R=f_ja-\sum_i(u_i f_jb_i+v_i f_jz_i)
                                                        \tag{136.LGF14}
\]
is closed. Indeed \(d(u_i f_jb_i)=f_jx_i f_jb_i-
N_i u_i f_jz_i\), while \(d(v_i f_jz_i)=
(N_i u_i-f_jy_i)f_jz_i\); their sum is \(f_j(da)\).

Choose \(dw_j=R_j\) with the \(w_j\) having common graph
restrictions. Here is the required existence argument. The
\(R_j|_G\) are identical and closed. Graph third cohomology
vanishes, so choose a common graph primitive. Target third
cohomology supplies arbitrary simplex primitives. Their graph
differences from the prescribed primitive are closed two-forms,
which are exact by (136.LGF1). Extend their one-form primitives
using the fixed row bound and subtract their differentials.
This makes the graph restrictions literally common.

Set
\[
 A_j=\sum_i(u_i f_jt_i+v_i f_jc_i)+w_j,\qquad
 B_j=\sum_i\ell_i v_i f_jb_i+Lw_j.
                                                        \tag{136.LGF15}
\]
All \(A_j\) have common graph restrictions. Direct differentiation
using (136.LGF9), (136.LGF13), and (136.LGF14) gives
\[
\begin{aligned}
 dA_j&=f_jX,& dB_j&=f_jh,\\
 LA_j-f_j\left(\sum_i\ell_i y_i t_i\right)
   &=d\left(\sum_i\ell_i v_i f_jt_i\right)+B_j.
\end{aligned}                                         \tag{136.LGF16}
\]
For example, the last identity uses
\(L u_i-\ell_i f_jy_i=\ell_i\,dv_i\);
the extra term from \(d(v_i f_jt_i)\) is exactly
\(-v_i f_jb_i\).

Old graph compatibility for the pair \(dH=Mh\) supplies
\[
 dA_{0j}=f_jh,\qquad dQ_{0j}=f_jH-MA_{0j},               \tag{136.LGF17}
\]
with both graph restrictions common. The closed two-form
\(B_j-A_{0j}\) is exact in the target simplex; choose a
one-form \(s_j\) with \(ds_j=B_j-A_{0j}\). Because the
right-hand sides have common graph restrictions,
\((s_j-s_0)|_G\) is a closed graph one-form. Put
\[
 Q_j=Q_{0j}-M s_j-M\sum_i\ell_i v_i f_jt_i.
                                                        \tag{136.LGF18}
\]
Then (136.LGF16)--(136.LGF17) give exactly
\[
\begin{aligned}
 dQ_j&=f_j\left(H+M\sum_i\ell_i y_i t_i\right)-ML A_j,\\
 \Delta_j([X])
   &=-\left[\frac{\operatorname{per}_G((s_j-s_0)|_G)}{L}\right]
          \quad\text{in }J^w .
\end{aligned}                                         \tag{136.LGF19}
\]
The source expression in the first line has differential \(MLX\)
by (136.LGF12). In the graph difference of (136.LGF18), all terms
except \(-M(s_j-s_0)\) cancel. Division by the order \(ML\)
in (136.LGF3) therefore gives the second line.

All elements defining \(R_j,B_j,A_{0j},s_j\) use only the
old maps and the old data (136.LGF10), not the new images of
\(c_i,t_i\). They can be kept fixed when those new images
are changed. This proves the stated independence from the
choice of the new extensions. The expression modulo \(S\)
is independent of the auxiliary choices because it equals the
intrinsic linking class. In particular, changing \(w_j\) by
an exact two-form with common graph differential changes the
displayed period only by a multiple of \(L\).
\end{proof}

\paragraph{The precise unresolved condition.}
The closed forms \(B_j-A_{0j}\) in (136.LGF17)--(136.LGF19)
are constructed using target primitives of old source data.
They are not generally the images of one old source two-cycle.
Consequently the torsion-to-torsion-free graph-period argument
for \(H^2(A_2)\) cannot be applied to their primitive differences.
The formula itself neither supplies a choice making this period
vanish nor proves a preservation theorem. The actual nonzero
example in ARCH-136-Q2-Q3-COMPATIBILITY-OBSTRUCTION shows that
this missing condition is substantive, even with the literal common
first skeleton. It remains an obstruction to retaining those old
maps, not a counterexample to the existence of a common source.

When the old linking map is zero, the universal linking map on
the arity-one part of the new \(H^3\) factors through the
\(\ker\Phi\) quotient of ARCH-136-Q3-LOW-FILTRATION.
Formula (136.LGF19) computes its finite homogeneous representatives
in the stated individual-weight guards. It does not control the
separate coefficient-dependent arity-two quotient.

\paragraph{Verification status.}
The author checked every displayed signed differential and completed
a whole saved-proof reading. The single-cell mixed formula and the
unequal-order Tor formula previously received direct root and cote15
checks in the research note. The root and cote15 each completed a whole saved-proof reading of
this finite-family LaTeX version and reported PASS. Their actual events
are recorded separately; the root also contributed the earlier discussion.
Neither statement is counted as a reduction of the unresolved main
common-source problem.
